% 数学ノート集・共通テンプレート v2（積分論2026年度第1部の版面を参考）
% A5判。章内で定義・定理・例・注意等の番号を共通化。
% コンパイル： platex -interaction=nonstopmode math_note_template_v2.tex (2回)
%              dvipdfmx math_note_template_v2.dvi
\documentclass[a5paper,10pt,onecolumn]{jsarticle}
\usepackage[dvipdfmx,top=22mm,bottom=20mm,inner=19mm,outer=18mm,headheight=13pt,headsep=6mm]{geometry}
\usepackage{amsmath,amssymb,amsthm,mathtools}
\usepackage{enumitem,fancyhdr,titlesec}
\usepackage[dvipdfmx,bookmarks=true,colorlinks=false]{hyperref}
\usepackage{pxjahyper}

% ----- 差し替える文書情報 -----
\newcommand{\NoteTitle}{Pell方程式と二次体の単数}
\newcommand{\NoteSeries}{数学ノート集}

% ----- 文字・組版 -----
% 参考PDFの質素な明朝＋Computer Modernの印象を引き継ぎながら，
% A5判で窮屈にならないよう本文の行間を少し確保する．
\linespread{1.13}
\setlength{\parindent}{1zw}
\setlength{\parskip}{0pt}
\setlength{\abovedisplayskip}{6pt plus 2pt minus 2pt}
\setlength{\belowdisplayskip}{6pt plus 2pt minus 2pt}
\setlength{\abovedisplayshortskip}{3pt plus 2pt minus 1pt}
\setlength{\belowdisplayshortskip}{4pt plus 2pt minus 1pt}
\setlist[enumerate]{topsep=3pt,itemsep=1pt,parsep=0pt,leftmargin=2.4zw}
\titleformat{\section}{\normalfont\normalsize\bfseries}{\thesection}{0.8zw}{}
\titlespacing*{\section}{0pt}{2.0ex plus .4ex}{.9ex}
\titleformat{\subsection}{\normalfont\normalsize\bfseries}{\thesubsection}{0.7zw}{}
\titlespacing*{\subsection}{0pt}{1.5ex plus .3ex}{.6ex}
\raggedbottom

% ----- 柱とページ番号：ページ上部、罫線なし -----
\fancypagestyle{mathnote}{%
  \fancyhf{}%
  \fancyhead[L]{\footnotesize \NoteSeries\ （\NoteTitle）}%
  \fancyhead[R]{\footnotesize \thepage}%
  \renewcommand{\headrulewidth}{0pt}%
  \renewcommand{\footrulewidth}{0pt}%
}
\pagestyle{mathnote}

% ----- 節内通し番号の定義・定理群 -----
\newtheoremstyle{mathnotetheorem}
 {4.5pt}{4.5pt}{\normalfont}{}%
 {\bfseries}{}{0.65zw}%
 {\thmname{#1}\thmnumber{ #2}\thmnote{（#3）}}
\theoremstyle{mathnotetheorem}
\newtheorem{theorem}{定理}[section]
\newtheorem{definition}[theorem]{定義}
\newtheorem{lemma}[theorem]{補題}
\newtheorem{proposition}[theorem]{命題}
\newtheorem{corollary}[theorem]{系}
\newtheorem{example}[theorem]{例}
\newtheorem{remark}[theorem]{注意}
\renewcommand{\proofname}{証明}
\renewcommand{\qedsymbol}{$\square$}
\makeatletter
\renewenvironment{proof}[1][\proofname]{%
  \par\pushQED{\qed}\normalfont\topsep4\p@\@plus2\p@\relax
  \trivlist\item[\hskip\labelsep\bfseries #1]\ignorespaces
}{\popQED\endtrivlist\@endpefalse}
\makeatother

% ----- 使うノートで適宜拡張する記号 -----
\newcommand{\R}{\mathbb{R}}
\newcommand{\norm}[1]{\lVert #1\rVert}
\newcommand{\abs}[1]{\lvert #1\rvert}

\begin{document}
\section*{読了後にできること}
本稿を読むための前提は，Pell方程式，二次体の単数とノルム．
本稿の具体的な到達課題は，Pell方程式の解をノルム1の単数として構成する．

\section{Pell方程式と二次環}

$D>1$ を平方数でない正整数とする．
\begin{definition}
方程式
\[
 x^2-Dy^2=1,\qquad x,y\in\mathbb Z
\]
を\emph{Pell方程式}という．
$R_D=\mathbb Z[\sqrt D]
 =\{a+b\sqrt D:a,b\in\mathbb Z\}$
と置き，
\[
 \overline{a+b\sqrt D}=a-b\sqrt D,\qquad
 N(a+b\sqrt D)=a^2-Db^2
\]
を共役およびノルムという．
\end{definition}

\begin{lemma}\label{lem:norm}
共役は $R_D$ の環自己同型であり，
任意の $\alpha,\beta\in R_D$ に対して
\[
 N(\alpha\beta)=N(\alpha)N(\beta)
\]
が成り立つ．
また $\alpha\in R_D$ が単元であることと
$N(\alpha)=\pm1$ は同値である．
\end{lemma}
\begin{proof}
共役は和と積を保つことを直接計算で確認でき，
共役を二回行うと恒等写像に戻る．
$N(\alpha)=\alpha\bar\alpha$ なので
\[
 N(\alpha\beta)
 =\alpha\beta\overline{\alpha\beta}
 =\alpha\bar\alpha\,\beta\bar\beta
 =N(\alpha)N(\beta).
\]
$\alpha$ が単元ならある $\beta\in R_D$ で
$\alpha\beta=1$ となるから
$N(\alpha)N(\beta)=1$ である．
両ノルムは整数なので $N(\alpha)=\pm1$ となる．
逆に $N(\alpha)=\pm1$ なら
$\alpha^{-1}=\bar\alpha/N(\alpha)=\pm\bar\alpha\in R_D$
であるから単元である．
\end{proof}

\begin{corollary}
Pell方程式の整数解 $(x,y)$ と，
ノルム1の単元 $x+y\sqrt D\in R_D$ は一対一に対応する．
\end{corollary}
\begin{proof}
$N(x+y\sqrt D)=x^2-Dy^2$ なので，
方程式とノルム1の条件は同値である．
また $D$ は平方数でないため $1,\sqrt D$ は
$\mathbb Q$ 上一次独立であり，
$(x,y)$ の表示は一意である．
\end{proof}

\section{初等的な近似による単元の存在}

\begin{lemma}[Dirichletの近似補題]\label{lem:dirichlet}
$\alpha\in\mathbb R\setminus\mathbb Q$ とする．
任意の正整数 $Q$ に対して，
整数 $p,q$ で $1\le q\le Q$ かつ
\[
 |q\alpha-p|<\frac1Q
\]
を満たすものが存在する．
さらに $Q\to\infty$ に沿って
このような組を選べば，その分母 $q$ は
ある部分列上で無限大に発散する．
\end{lemma}
\begin{proof}
$0,\alpha,2\alpha,\dots,Q\alpha$ の小数部分
$Q+1$ 個を $[0,1)$ の長さ $1/Q$ の半開区間 $Q$ 個に
割り当てる．鳩の巣原理により，
二つの小数部分が同じ区間に入る．
その差に対応する整数 $q$ は $1\le q\le Q$ を満たし，
適切な整数 $p$ に対して
$|q\alpha-p|<1/Q$ となる．
もし選ばれた $q$ が有界な無限部分列を持つなら，
さらに同じ整数 $q$ が無限回現れる部分列を取れる．
しかし無理数 $\alpha$ に対して
$|q\alpha-p|$ は整数 $p$ を動かしても
$0$ より大きい最小値を持つため，
$1/Q\to0$ に反する．
従って $q$ は適切な部分列上で無限大に発散する．
\end{proof}

\begin{lemma}[有界なノルムを持つ無限列]\label{lem:bounded-norm}
$D>1$ を平方数でない正整数とする．
正整数 $q_n\to\infty$ と整数 $p_n>0$ の列で
\[
 0<|p_n^2-Dq_n^2|\le C_D
\]
を満たすものが存在する．
ここで $C_D$ は $D$ のみで決まる定数である．
\end{lemma}
\begin{proof}
補題\ref{lem:dirichlet}を $\alpha=\sqrt D$ に適用し，
$q_n\to\infty$ となる近似
$|q_n\sqrt D-p_n|<1/q_n$ を取る．
$q_n$ が十分大きいとき $p_n>0$ である．
このとき
\[
 |p_n^2-Dq_n^2|
 =|p_n-q_n\sqrt D|(p_n+q_n\sqrt D)
 <\frac1{q_n}\left(2q_n\sqrt D+\frac1{q_n}\right)
 \le2\sqrt D+1.
\]
左辺は整数であり，$D$ が平方数でないため零にはならない．
従って $C_D=2\sqrt D+1$ と取ればよい．
\end{proof}

\begin{theorem}[Pell方程式の非自明解の存在]\label{thm:pell-exists}
$D>1$ を平方数でない正整数とする．
このとき $x>1$，$y>0$ を満たすPell方程式の整数解
$(x,y)$ が存在する．
\end{theorem}
\begin{proof}
補題\ref{lem:bounded-norm}により，
$\alpha_n=p_n+q_n\sqrt D$ のノルムは
零でない有界な整数であり，
$\alpha_n\to\infty$ である．
従ってある零でない整数 $m$ に対して
\[
 N(\alpha_n)=m
\]
となる $\alpha_n$ が無限個存在する．
さらに係数 $(p_n,q_n)$ を法 $|m|$ で考えれば
可能な剰余類は有限個なので，
同じ剰余類を持ち，かつ異なる正の実数値を持つ
二つの元
$\alpha=a+b\sqrt D$，
$\beta=c+d\sqrt D$ を取れる．
必要なら順序を入れ替えて $\alpha>\beta>0$ とする．
これらは
\[
 a\equiv c,\quad b\equiv d\pmod{|m|}
\]
および $N(\alpha)=N(\beta)=m$ を満たす．
ここで
\[
 \varepsilon=\frac{\alpha}{\beta}
 =\frac{\alpha\bar\beta}{m}
 =\frac{(ac-Db d)+(bc-ad)\sqrt D}{m}
\]
と置く．
剰余条件から
$ac-Db d\equiv c^2-Dd^2=m\equiv0\pmod{|m|}$
および
$bc-ad\equiv dc-cd=0\pmod{|m|}$
である．従って
$\varepsilon\in\mathbb Z[\sqrt D]$ である．
ノルムの乗法性から
$N(\varepsilon)=N(\alpha)/N(\beta)=1$ であり，
また $\alpha>\beta>0$ だから $\varepsilon>1$ である．
$\varepsilon=x+y\sqrt D$ と書くと
$N(\varepsilon)=1$ から $x^2-Dy^2=1$ となる．
さらに共役は
$\bar\varepsilon=1/\varepsilon\in(0,1)$ であるから
\[
 x=\frac{\varepsilon+\bar\varepsilon}2>1,\qquad
 y=\frac{\varepsilon-\bar\varepsilon}{2\sqrt D}>0.
\]
よって求める非自明な解を得る．
\end{proof}

\begin{remark}
定理\ref{thm:pell-exists}は連分数展開を用いていない．
無理数の有理近似から有界なノルムを持つ元を無限に作り，
同じノルムと合同類を持つ二元の商を取ることによって
ノルム1の単元を構成している．
\end{remark}

\section{基本単元と解の分類}

\begin{lemma}[正の単元の離散性]\label{lem:discrete}
$D>1$ を平方数でない整数とする．
任意の $M>1$ に対し
\[
 \{u\in R_D:N(u)=1,\ 1\le u\le M\}
\]
は有限集合である．
\end{lemma}
\begin{proof}
$u=x+y\sqrt D$ がノルム1を満たすなら
$\bar u=1/u$ である．
従って
\[
 x=\frac{u+1/u}2,\qquad
 y=\frac{u-1/u}{2\sqrt D}
\]
となる．$1\le u\le M$ のとき $x,y$ は有界な実数区間に
属し，しかも整数である．従って可能な組 $(x,y)$ は有限個であり，
その集合も有限である．
\end{proof}

\begin{theorem}[基本単元とPell方程式の全解]\label{thm:fundamental}
$D>1$ を平方数でない正整数とする．
ノルム1の単元 $u>1$ の中に最小のもの
\[
 \varepsilon_D=a+b\sqrt D>1
\]
が存在する．
さらにノルム1の単元全体は
\[
 \{u\in R_D:N(u)=1\}
 =\{\pm\varepsilon_D^n:n\in\mathbb Z\}
\]
と一致する．
従って全てのPell方程式の整数解は
$\pm\varepsilon_D^n$ の係数として得られる．
\end{theorem}
\begin{proof}
定理\ref{thm:pell-exists}により
ノルム1の単元で $u>1$ となるものが存在する．
その一つを $u_0>1$ とする．
補題\ref{lem:discrete}によって
$1<u\le u_0$ を満たす単元は有限個なので，
その中に最小のもの $\varepsilon_D$ が存在する．
これは全ての $u>1$ の中で最小である．

任意のノルム1の単元 $u>0$ を取る．
$\varepsilon_D>1$ なので，
ある整数 $n$ が存在して
\[
 \varepsilon_D^n\le u<\varepsilon_D^{n+1}
\]
となる．$v=u\varepsilon_D^{-n}$ と置くと
$v\in R_D$ はノルム1の単元であり，
$1\le v<\varepsilon_D$ である．
最小性から $v=1$ でなければならない．
従って $u=\varepsilon_D^n$ である．
任意の負のノルム1単元 $u$ については
$-u>0$ であり，$N(-u)=N(u)=1$ だから
$-u=\varepsilon_D^n$ と書ける．
従って全ての元は $\pm\varepsilon_D^n$ となる．
逆にこれらはノルムの乗法性から全てノルム1である．
\end{proof}

\begin{corollary}[Pell解の漸化式]\label{cor:recurrence}
$\varepsilon_D=a+b\sqrt D$ とし，
\[
 \varepsilon_D^n=x_n+y_n\sqrt D,\qquad n\ge0
\]
と置く．このとき
\begin{align*}
 x_0&=1,\quad y_0=0,\\
 x_{n+1}&=ax_n+Db\,y_n,\\
 y_{n+1}&=bx_n+ay_n.
\end{align*}
さらに $n\ge0$ の各 $(x_n,y_n)$ はPell方程式の解である．
\end{corollary}
\begin{proof}
$\varepsilon_D^{n+1}
 =(a+b\sqrt D)(x_n+y_n\sqrt D)$ を展開し，
$1$ と $\sqrt D$ の係数を比較すれば漸化式を得る．
$N(\varepsilon_D^n)=1$ なので
$x_n^2-Dy_n^2=1$ である．
\end{proof}

\section{ノルムが負の単元}

\begin{proposition}[負のPell方程式]\label{prop:negative}
$D>1$ を平方数でない正整数とする．
$R_D$ にノルム $-1$ の単元が存在するなら，
ノルム1の基本単元 $\varepsilon_D$ は
ノルム $-1$ の単元 $\eta$ の平方
$\varepsilon_D=\eta^2$ と書ける．
このときノルム $-1$ の単元全体は
$\{\pm\eta\,\varepsilon_D^n:n\in\mathbb Z\}$ となる．
\end{proposition}
\begin{proof}
$N(\alpha)=-1$ となる単元 $\alpha$ を取る．
必要なら符号と逆元を使って
$\alpha>1$ としてよい．
$\alpha^2$ はノルム1で正の単元だから
定理\ref{thm:fundamental}より
$\alpha^2=\varepsilon_D^k$（$k\ge1$）と書ける．
もし $k=2j$ が偶数なら，
$(\alpha\varepsilon_D^{-j})^2=1$ となる．
左辺の元は正の実数なので
$\alpha\varepsilon_D^{-j}=1$ である．
しかし左辺のノルムは $-1$ であり矛盾する．
従って $k=2j+1$ は奇数である．
$\eta=\alpha\varepsilon_D^{-j}$ と置けば
$N(\eta)=-1$ かつ $\eta^2=\varepsilon_D$ となる．

任意のノルム $-1$ の単元 $\beta$ に対し
$\beta/\eta$ はノルム1の単元である．
定理\ref{thm:fundamental}から
$\beta/\eta=\pm\varepsilon_D^n$ と書けるので
$\beta=\pm\eta\varepsilon_D^n$ となる．
逆にこれらのノルムは全て $-1$ である．
\end{proof}

\section{具体例：$D=2$ と $D=7$}

\begin{example}[$D=2$]
方程式 $x^2-2y^2=1$ のノルム1の基本単元は
\[
 \varepsilon_2=3+2\sqrt2
\]
である．従って全解は
$x+y\sqrt2=\pm(3+2\sqrt2)^n$（$n\in\mathbb Z$）である．
また $1+\sqrt2$ はノルム $-1$ の単元であり，
$(1+\sqrt2)^2=3+2\sqrt2$ となる．
\end{example}
\begin{proof}
$N(3+2\sqrt2)=9-8=1$ である．
任意のノルム1の単元 $u=x+y\sqrt2>1$ は
$x>1$，$y>0$ である．
$y=1$ なら $x^2=3$ となって整数解を持たず，
$y=2$ なら $x=3$ が得られる．
$y$ が増大すると
$u=x+y\sqrt2$ も増大するので，
$3+2\sqrt2$ が最小のノルム1単元である．
$N(1+\sqrt2)=1-2=-1$ と平方の恒等式は
直接計算で従う．
\end{proof}

\begin{example}[$D=7$]
方程式 $x^2-7y^2=1$ のノルム1の基本単元は
\[
 \varepsilon_7=8+3\sqrt7
\]
である．また負のPell方程式
$x^2-7y^2=-1$ は整数解を持たない．
\end{example}
\begin{proof}
$8^2-7\cdot3^2=64-63=1$ である．
$y=1$ なら $x^2=8$，$y=2$ なら $x^2=29$ なので
いずれも整数解はない．
$y=3$ で初めて正の解 $x=8$ を得る．
$u=x+y\sqrt7>1$ のとき $x,y$ は正であり，
$y$ に関して $u$ は増加するため
$8+3\sqrt7$ が最小の正のノルム1単元である．
負のPell方程式が解を持つなら
$x^2\equiv-1\equiv6\pmod7$ となるが，
法7の平方剰余は $0,1,2,4$ のみなので不可能である．
\end{proof}

\section{二次整数環との比較}

以下，$D>1$ は平方因子を持たない整数とする．

\begin{theorem}[実二次体の整数環]\label{thm:integers}
$K=\mathbb Q(\sqrt D)$ の整数環 $\mathcal O_K$ は
\[
 \mathcal O_K=
 \begin{cases}
 \mathbb Z[\sqrt D]&(D\equiv2,3\pmod4),\\[2mm]
 \mathbb Z[(1+\sqrt D)/2]&(D\equiv1\pmod4)
 \end{cases}
\]
に等しい．
\end{theorem}
\begin{proof}
$\alpha=a+b\sqrt D$（$a,b\in\mathbb Q$）が
代数的整数なら，そのトレース $2a$ とノルム
$a^2-Db^2$ は整数である．
これは代数的整数の共役も代数的整数であり，
和と積が有理数である代数的整数は整数であることから従う．
$u=2a\in\mathbb Z$ と置く．
最小多項式の判別式は
\[
 u^2-4(a^2-Db^2)=4Db^2.
\]
従って $4Db^2$ は整数の平方因子を含めた整数である．
$b\in\mathbb Q$ とし $2b=v/w$ を既約分数で表せば，
$D v^2/w^2\in\mathbb Z$ となる．
$D$ は平方因子を持たないので，
素数ごとの指数を比較すれば $w=1$ でなければならない．
従って $v=2b\in\mathbb Z$ であり，
\[
 \alpha=\frac{u+v\sqrt D}{2}
\]
と書ける．
ノルムは
$(u^2-Dv^2)/4\in\mathbb Z$ なので
$u^2-Dv^2\equiv0\pmod4$ である．
$D\equiv2,3\pmod4$ のとき
$u,v$ がともに偶数でなければこの合同式は成り立たない．
従って $\alpha\in\mathbb Z[\sqrt D]$ となる．
$D\equiv1\pmod4$ のときは $u,v$ が同じ偶奇であることと
合同式が同値であり，
$\alpha\in\mathbb Z[(1+\sqrt D)/2]$ となる．
逆に表示された環の元は，
それぞれ $\sqrt D$ または $(1+\sqrt D)/2$ が
モニックな整数係数多項式の根であるため
代数的整数である．
よって等号を得る．
\end{proof}

\begin{example}[$D=5$]
$\mathbb Q(\sqrt5)$ の整数環は
$\mathbb Z[(1+\sqrt5)/2]$ であり，
\[
 \eta=\frac{1+\sqrt5}{2}
\]
は整数環のノルム $-1$ の単元である．
しかし $\eta\notin\mathbb Z[\sqrt5]$ である．
一方，$x^2-5y^2=1$ の基本単元は
$9+4\sqrt5$ である．
\end{example}
\begin{proof}
$D=5\equiv1\pmod4$ なので整数環の表示は
定理\ref{thm:integers}による．
$N(\eta)=(1-5)/4=-1$ であり，
$\eta$ の $\sqrt5$ の係数は $1/2$ だから
$\eta\notin\mathbb Z[\sqrt5]$ である．
$9^2-5\cdot4^2=81-80=1$ であり，
$y=1,2,3$ に対して
$x^2=6,21,46$ はいずれも平方数でない．
従って $9+4\sqrt5$ が
$\mathbb Z[\sqrt5]$ における最小のノルム1単元である．
\end{proof}

\begin{remark}
Pell方程式の全解は環 $\mathbb Z[\sqrt D]$ の
ノルム1単元によって記述される．
二次体の整数環がそれより大きい場合，
整数環の単元群とPell方程式の解群を
区別する必要がある．
\end{remark}

\section{到達課題の完全解答}
\noindent\textbf{Pell方程式の解を単数から生成する．}
\begin{proposition}
任意の $n\ge0$ に対し
\[
x_n+y_n\sqrt2=(3+2\sqrt2)^n
\]
で定める整数 $x_n,y_n$ は
$x_n^2-2y_n^2=1$ を満たす．
\end{proposition}
\begin{proof}
ノルム
$N(a+b\sqrt2)=a^2-2b^2$ は乗法的であり，
$N(3+2\sqrt2)=9-8=1$．
従って $N((3+2\sqrt2)^n)=1$ である．
二項定理により係数 $x_n,y_n$ は整数である．
\end{proof}
\begin{proposition}
上の列は
$x_{n+1}=3x_n+4y_n$，
$y_{n+1}=2x_n+3y_n$ を満たす．
\end{proposition}
\begin{proof}
$(x_n+y_n\sqrt2)(3+2\sqrt2)$ を展開して
有理部分と $\sqrt2$ の係数を比較する．
\end{proof}
\noindent\textbf{到達確認．}
$n=2$ の解を求めよ．
\par\noindent\textbf{解答．}
$(3+2\sqrt2)^2=17+12\sqrt2$ より
$(x_2,y_2)=(17,12)$．
実際 $17^2-2\cdot12^2=289-288=1$．
\end{document}